\documentclass{beamer}

% Theme selection
\usetheme{Madrid}
\usecolortheme{default}

% Packages
\usepackage[utf8]{inputenc}
\usepackage[T1]{fontenc}
\usepackage{amsmath}
\usepackage{amssymb}
\usepackage{graphicx}

% Title information
\title[Quantum Bayesian Game Framework]{Deep Dive: Dari Classical Bayesian Game ke Quantum EWL Framework}
\subtitle{Transisi Teoritis dan Arsitektur Sistem}
\author{DreamyTA Team}
\date{\today}

\begin{document}

% Slide 1: Title
\begin{frame}
    \titlepage
\end{frame}

% Slide 2: Daftar Isi
\begin{frame}{Daftar Isi}
    \tableofcontents
\end{frame}



\section{Landasan Teori: Classical Bayesian Game Theory}

% Slide 4: Konsep Dasar
\begin{frame}{Landasan Teori: Classical Bayesian Game Theory}
    \begin{block}{Definisi dan Konsep Dasar}
        \begin{itemize}
            \item \textbf{Masalah:} Di pasar nyata, investor tidak memiliki informasi sempurna (\textit{incomplete information}).
            \item \textbf{Bayesian Game:} Menangani situasi di mana pemain tidak mengetahui "tipe" dari lawan ("Nature").
            \item \textbf{Komponen:}
            \begin{itemize}
                \item \textbf{Pemain:} Investor/Algoritma Portofolio.
                \item \textbf{Nature:} Pasar Saham.
                \item \textbf{Tipe:} Kondisi Pasar (Bullish/Bearish).
                \item \textbf{Prior Belief:} Probabilitas tipe pasar (misal: 60\% Bull).
            \end{itemize}
        \end{itemize}
    \end{block}
\end{frame}

% Slide 5: Transformasi Harsanyi
\begin{frame}{Transformasi Harsanyi}
    Mengubah permainan informasi tidak lengkap menjadi permainan informasi tidak sempurna namun lengkap.
    
    \begin{enumerate}
        \item \textbf{Nature Move:} Alam memilih tipe pasar (Bull/Bear) sesuai probabilitas.
        \item \textbf{Player Move:} Investor memilih strategi alokasi aset.
        \item \textbf{Payoff:} Bergantung pada strategi dan tipe pasar.
    \end{enumerate}

    \vspace{0.5cm}
    \textbf{Ekspektasi Payoff:}
    $$ \mathbb{E}[\text{Payoff}] = \sum_{t \in \text{Types}} P(t) \cdot u(\text{Strategi}, t) $$
\end{frame}

\section{Landasan Teori: Skema EWL}

% Slide 6: Mengapa Perlu Kuantum?
\begin{frame}{Landasan Teori: Skema EWL}
    \begin{itemize}
        \item \textbf{Superposisi:} Strategi bukan biner, tapi kombinasi linier basis strategi.
        \item \textbf{Entanglement:} Korelasi non-lokal antar aset yang lebih kuat dari korelasi klasik.
        \item \textbf{Interferensi:} Memungkinkan solusi Pareto Optimal yang tidak terjangkau secara klasik.
    \end{itemize}
\end{frame}

% Slide 7: Protokol Standar EWL
\begin{frame}{Protokol Standar EWL}
    Dipetakan ke dalam 3 tahap sirkuit kuantum:
    
    \begin{enumerate}
        \item \textbf{Entangling Operator ($\hat{J}$):}
        $$ |\psi_{ent}\rangle = \hat{J} |00\dots0\rangle $$
        Biasanya $\hat{J} = \exp(i \frac{\gamma}{2} \sigma_x^{\otimes n}) $$ untuk korelasi maksimal.
        
        \item \textbf{Strategy Operator ($\hat{U}$):}
        Rotasi parameter ($\theta, \phi$) merepresentasikan bobot alokasi.
        $$ |\psi_{strat}\rangle = (\hat{U}_1 \otimes \hat{U}_2 \otimes \dots) |\psi_{ent}\rangle $$
        
        \item \textbf{Disentangling Operator ($\hat{J}^\dagger$):}
        $$ |\psi_{fin}\rangle = \hat{J}^\dagger |\psi_{strat}\rangle $$
    \end{enumerate}
\end{frame}

\section{Formulasi Matematis \& Hamiltonian}

% Slide 8: Fungsi Objektif (Markowitz)
\begin{frame}{1. Fungsi Objektif: Model Markowitz}
    Tujuan utama adalah meminimalkan risiko (varians) sekaligus memaksimalkan return.
    
    \begin{block}{Persamaan Dasar}
    $$ \text{Minimize } \quad C(x) = w \sum_{i,j} \sigma_{ij} x_i x_j - (1-w) \sum_i \mu_i x_i $$
    \end{block}
    
    \begin{itemize}
        \item $x_i \in \{0, 1\}$: Variabel keputusan biner (Aset dipilih/tidak).
        \item $\sigma_{ij}$: Kovarians antar aset $i$ dan $j$ (Risiko).
        \item $\mu_i$: Ekspektasi return aset $i$.
        \item $w$: Parameter penghindaran risiko (Risk Aversion, $0 \le w \le 1$).
    \end{itemize}
\end{frame}

% Slide 9: Constraints & Penalty
\begin{frame}{2. Kendala \& Penalti (Penalty Term)}
    VQE meminimalkan Hamiltonian tanpa kendala, sehingga kendala harus diubah menjadi "Soft Constraint" berupa penalti energi.
    
    \begin{block}{Kendala Kardinalitas}
    Kita ingin memilih tepat $B$ aset:
    $$ \sum_{i=1}^N x_i = B $$
    \end{block}
    
    \begin{alertblock}{Penalty Term}
    $$ P \left( \sum_{i=1}^N x_i - B \right)^2 $$
    Dimana $P$ adalah koefisien penalti yang besar.
    \end{alertblock}
\end{frame}

% Slide 10: Mapping ke Ising Hamiltonian
\begin{frame}{3. Mapping ke Ising Hamiltonian}
    Komputer kuantum bekerja dengan qubit (spin), bukan bit biner $0/1$.
    
    \textbf{Transformasi Variabel:}
    $$ x_i = \frac{1 - Z_i}{2} \quad \text{dimana } Z_i \in \{1, -1\} $$
    
    \vspace{0.3cm}
    \textbf{Hamiltonian Akhir (QUBO to Ising):}
    Substitusi $x_i$ ke dalam fungsi biaya menghasilkan operator Pauli-Z:
    \begin{equation*}
        H = \underbrace{\sum_i h_i Z_i}_{\text{Field Eksternal}} + \underbrace{\sum_{i<j} J_{ij} Z_i Z_j}_{\text{Interaksi (Coupling)}} + C
    \end{equation*}
    Inilah Hamiltonian yang akan dicari Ground State-nya oleh sirkuit EWL.
\end{frame}

\section{Arsitektur Sistem (Fase 1 \& 2)}

% Slide 11: Arsitektur Sistem Revised
\begin{frame}{Arsitektur Sistem: Input \& Pembentukan Model}
    \begin{center}
        \textbf{Fokus: Konstruksi Model Teoritis}
    \end{center}
    
    \begin{block}{Fase 1: Input \& Teori Dasar}
        \begin{itemize}
            \item \textbf{Data Historis} $\xrightarrow{\text{Statistik}}$ Risk ($\sigma_{ij}$) \& Return ($\mu_i$)
            \item \textbf{Game Theory} $\xrightarrow{\text{Incomplete Info}}$ Bayesian Game
            \item \textbf{Quantum Mechanics} $\xrightarrow{\text{Superposition}}$ EWL Scheme
        \end{itemize}
    \end{block}
    
    \begin{block}{Fase 2: Pembentukan Hamiltonian}
        \begin{itemize}
            \item \textbf{Langkah 1:} Tentukan Matriks Payoff (Basis Markowitz).
            \item \textbf{Langkah 2:} Terapkan Penalti $P(\sum x - B)^2$.
            \item \textbf{Langkah 3:} Transformasi Biner ke Spin ($x \to Z$).
            \item \textbf{Output:} \textbf{Bayesian Weighted Hamiltonian} ($H_{total}$).
        \end{itemize}
    \end{block}
\end{frame}

\section{Kesimpulan}

% Slide 12: Kesimpulan
\begin{frame}{Kesimpulan Formulasi}
    \begin{enumerate}
        \item \textbf{Integrasi Hibrida:} 
        Kita berhasil memetakan masalah optimasi portofolio klasik (Markowitz) ke dalam domain fisika kuantum (Ising Model).
        
        \item \textbf{Peran Game Theory:}
        Fungsi bobot ($w$) dan parameter penalti ($P$) dapat diadaptasi secara dinamis berdasarkan outcome dari Bayesian Game (Kondisi Pasar).
        
        \item \textbf{Siap untuk Eksekusi:}
        Hamiltonian $H$ yang telah dikonstruksi siap menjadi input bagi algoritma VQE, dimana sirkuit EWL akan bertindak sebagai \textit{Ansatz} pencari solusi.
    \end{enumerate}
\end{frame}

\end{document}
